% Figure: the shooting method for a two-point BVP. Several trial
% trajectories with different initial slopes; the correct slope hits
% the target boundary value at x = L.
\begin{figure}[htbp]
\centering
\begin{tikzpicture}[
  axis/.style={draw=engblack, -{Stealth}, thick},
  miss/.style={draw=enggray, thick},
  hit/.style={draw=engred, very thick},
  target/.style={draw=engblue, fill=engblue, circle, inner sep=1.5pt},
  label/.style={font=\small},
]
% Axes: t from 0 to 5 (independent), y from 0 to 4 (state)
\draw[axis] (-0.2,0) -- (5.6,0) node[anchor=west, label] {$t$};
\draw[axis] (0,-0.2) -- (0,4.4) node[anchor=south, label] {$y$};
% Boundary markers
\node[font=\scriptsize, below] at (0,-0.12) {$0$};
\node[font=\scriptsize, below] at (5,-0.12) {$L$};
\draw[densely dotted, enggray] (5,0) -- (5,4.2);
% Left boundary value y(0) = 1
\draw[engblack, fill=engblack] (0,1) circle (1.4pt);
\node[font=\scriptsize, left] at (-0.1,1) {$y_{0}$};
% Target right boundary value y(L) = 3
\node[target] at (5,3) {};
\node[font=\scriptsize, right, engblue] at (5.1,3) {$y_{L}$ (target)};
% Trial 1: slope too low, undershoots
\draw[miss] plot[smooth, domain=0:5, samples=60]
  (\x, {1 + 0.15*\x + 0.02*\x*\x});
\node[font=\scriptsize, enggray, anchor=west] at (5.05,1.25) {$s_{1}$ low};
% Trial 2: slope too high, overshoots
\draw[miss] plot[smooth, domain=0:5, samples=60]
  (\x, {1 + 0.75*\x - 0.01*\x*\x});
\node[font=\scriptsize, enggray, anchor=west] at (5.05,4.15) {$s_{2}$ high};
% Correct trajectory: hits the target
\draw[hit] plot[smooth, domain=0:5, samples=60]
  (\x, {1 + 0.48*\x - 0.004*\x*\x});
\node[font=\scriptsize, engred, anchor=south] at (2.5,2.5) {$s^{\ast}$ (solves BVP)};
\end{tikzpicture}
\caption[The shooting method for a boundary-value problem]{The
shooting method converts a two-point boundary-value problem into a
sequence of initial-value problems. The left boundary condition
$y(0) = y_{0}$ is fixed; the unknown initial slope $s = y'(0)$ is the
shooting parameter. Each trial slope produces one IVP trajectory
(grey). The residual $\phi(s) = y(L; s) - y_{L}$ measures how far the
trajectory misses the right boundary value; a root-finder (bisection
or Newton on $\phi$) adjusts $s$ until $\phi(s^{\ast}) = 0$ and the
trajectory (red) lands on the target. For a linear BVP the residual
is affine in $s$, so two shots and one linear interpolation give the
exact answer.}
\label{fig:vol02:ch07:shooting}
\end{figure}
